\PassOptionsToPackage{dvipsnames,svgnames,table}{xcolor}
\documentclass[11pt]{article}
\usepackage[a4paper,margin=22mm]{geometry}
\usepackage[no-math]{fontspec}
\setmainfont{DejaVu Serif}
\usepackage{amsmath,amssymb,amsthm,mathtools,graphicx,xcolor,hyperref,xurl,xspace,ifthen,enumerate}
\usepackage[shortlabels]{enumitem}
\setlength{\emergencystretch}{4em}
\AtBeginDocument{\special{pdf:mapline pzdr Dingbats <uzdr.pfb}}
\SetMathAlphabet{\mathbf}{normal}{OT1}{cmr}{bx}{n}
\SetMathAlphabet{\mathbf}{bold}{OT1}{cmr}{bx}{n}
\usepackage{thm-restate}
\usepackage{cleveref}
\newtheorem{theo}{Theorem}[section]
\newtheorem{lemm}[theo]{Lemma}
\newtheorem{thm}[theo]{Theorem}
\newtheorem{proposition}[theo]{Proposition}
\newtheorem{corollary}[theo]{Corollary}
\newtheorem{definition}[theo]{Definition}
\newtheorem{theorem}[theo]{Theorem}
\newtheorem*{remas*}{Remarks}
\newtheorem{remas}[theo]{Remarks}
\newtheorem*{exam*}{Example}
\newtheorem{ques}[theo]{Question}
\providecommand{\og}{«\nobreakspace}
\providecommand{\fg}{\nobreakspace»}
\newtheorem{lem}[theo]{Lemma}
\newtheorem{lemma}[theo]{Lemma}
\newtheorem{prop}[theo]{Proposition}
\newtheorem{coro}[theo]{Corollary}
\newtheorem{corol}[theo]{Corollary}
\newtheorem*{conj*}{Conjecture}
\newtheorem{conj}[theo]{Conjecture}
\newtheorem{conjecture}[theo]{Conjecture}
\newtheorem{Proposition}[theo]{Proposition}
\newtheorem{theoreme}[theo]{Theorem}
\theoremstyle{definition}
\newtheorem{defi}[theo]{Definition}
\newtheorem{exam}[theo]{Example}
\newtheorem{exem}[theo]{Example}
\newtheorem{remark}[theo]{Remark}
\newtheorem{example}[theo]{Example}
\newtheorem{exams}[theo]{Examples}
\newtheorem{rema}[theo]{Remark}
\newtheorem{nota}[theo]{Notation}
\newtheorem{notation}[theo]{Notation}
\newtheorem*{notas}{Notations}
\newtheorem{result}[theo]{Result}
\newtheorem{question}[theo]{Question}
\newtheorem{prob}[theo]{Problem}
\newtheorem*{rema*}{Remark}
\newtheorem*{defi*}{Definition}
\newtheorem*{theo*}{Theorem}
\newtheorem*{lemm*}{Lemma}
\newtheorem*{prop*}{Proposition}
\newcommand{\enoncename}{}
\newtheorem{corpusenonce}[theo]{\enoncename}
\newenvironment{enonce}[2][]{\renewcommand{\enoncename}{#2}\begin{corpusenonce}}{\end{corpusenonce}}
\newenvironment{enonce*}[2][]{\par\medskip\noindent\textbf{#2.}\itshape}{\par\medskip}
\newenvironment{proofc}[1][\proofname]{\begin{proof}[#1]}{\end{proof}}
\newcommand{\Mc}{\mathcal{M}}
\newcommand{\pointir}{.\ ---\ }
\newcommand{\equalenv}[2]{\expandafter\let\csname#1\expandafter\endcsname\csname#2\endcsname\expandafter\let\csname end#1\expandafter\endcsname\csname end#2\endcsname}
\providecommand{\diagbox}[3][]{\begin{tikzpicture}[x=1em,y=1em,baseline=(current bounding box.center)]\useasboundingbox(0,0) rectangle(3,1.4);\draw(0,1.4)--(3,0);\node at(.5,.3){#2};\node at(2.5,1.1){#3};\end{tikzpicture}}
\providecommand{\eg}{e.g.\ }
\providecommand{\ie}{i.e.\ }
\providecommand{\cf}{cf.\ }
\providecommand{\longthanks}[1]{\par\smallskip\noindent #1\par\smallskip}
\providecommand{\qbinom}[2]{\genfrac{[}{]}{0pt}{}{#1}{#2}}

\providecommand{\mathof}[1]{\mathrm{#1}}

\numberwithin{equation}{section}


\def\endappendix{}



\usepackage{cleveref}
\usepackage{thmtools}
\usepackage{mathrsfs}
\usepackage{tikz}
\usepackage{mathtools}
\newfontfamily\CorpusFiniteLMMath[Path=/usr/share/texmf/fonts/opentype/public/lm-math/]{latinmodern-math.otf}
\DeclareRobustCommand{\llbracket}{\mathopen{\text{\CorpusFiniteLMMath\char"27E6}}}
\DeclareRobustCommand{\rrbracket}{\mathclose{\text{\CorpusFiniteLMMath\char"27E7}}}


\allowdisplaybreaks











\DeclareMathOperator{\Set}{Set}
\DeclareMathOperator{\Id}{Id}

\newcommand{\statement}[1]{
  #1\enspace\ignorespaces
}



\newcommand*{\ldblbrace}{\{\mskip-5mu\{}
\newcommand*{\rdblbrace}{\}\mskip-5mu\}}
\DeclareRobustCommand{\stirlingi}{\genfrac[]{0pt}{}}
\DeclareRobustCommand{\stirlingii}{\genfrac\{\}{0pt}{}}
\newcommand{\widelangle}{{\mathpalette\scaledlangle\relax}}
\newcommand{\widerangle}{{\mathpalette\scaledrangle\relax}}
\newcommand{\scaledlangle}[2]{\scalebox{1.5}[1]{$#1\langle$}}
\newcommand{\scaledrangle}[2]{\scalebox{1.5}[1]{$#1\rangle$}}
\newcommand{\K}{\mathbb{C}}
\newcommand{\Z}{\mathbb{Z}}
\newcommand{\N}{\mathbb{N}}
\newcommand{\bF}{\mathbf{F}}
\newcommand{\bG}{\mathbf{G}}
\newcommand{\bS}{\mathbf{S}}
\newcommand{\Vect}{\text{Vect}}
\newcommand{\FinSet}{\text{FinSet}}
\newcommand{\Mod}{\mathrm{Mod}}
\newcommand{\Fa}{\mathscr{F}}
\newcommand{\Inj}{\mathchoice{\operatorname{Inj}}{\operatorname{Inj}}{\mathrm{Inj}}{\mathrm{Inj}}}
\newcommand{\Fun}{\mathchoice{\operatorname{Fun}}{\operatorname{Fun}}{\mathrm{Fun}}{\mathrm{Fun}}}
\newcommand{\Sur}{\mathchoice{\operatorname{Sur}}{\operatorname{Sur}}{\mathrm{Sur}}{\mathrm{Sur}}}
\newcommand{\Bij}{\mathchoice{\operatorname{Bij}}{\operatorname{Bij}}{\mathrm{Bij}}{\mathrm{Bij}}}
\newcommand{\rmSur}{\mathrm{Sur}}
\newcommand{\rmBij}{\mathrm{Bij}}
\newcommand{\rmFun}{\mathrm{Fun}}
\newcommand{\rmInj}{\mathrm{Inj}}
\newcommand{\F}[1]{\mathscr{F}\{#1\}}
\newcommand{\FSur}[1]{\mathscr{F}_{\Sur}\left\{#1\right\}}
\newcommand{\FSurInv}[1]{\mathscr{F}_{\Sur}^{-1}\left\{#1\right\}}
\newcommand{\E}{\mathbb{E}}
\newcommand{\Sym}{\mathfrak{S}}
\newcommand{\sqop}{\mathbin{\square}}
\DeclareMathOperator{\ch}{ch}
\DeclareMathOperator{\Hom}{Hom}
\DeclareMathOperator{\Lyndon}{Lyndon}
\newcommand{\branchtree}[6]{
    \begin{scope}[shift={(#1, -#2)}]
    \draw[line width=0.5pt] (0, -0.05) -- (0, 0.3);
    \draw[line width=0.5pt] (0, -0.05) -- (-0.3, -0.3);
    \draw[line width=0.5pt] (0, -0.05) -- (0.3, -0.3);
    \draw[#3, fill=#3] (0, -0.05) circle [radius=1pt];
    \draw[#4, fill=#4] (0, 0.3) circle [radius=1pt];
    \draw[#5, fill=#5] (-0.3, -0.3) circle [radius=1pt];
    \draw[#6, fill=#6] (0.3, -0.3) circle [radius=1pt];
    \end{scope}
}
\newcommand{\squaretree}[6]{
    \begin{scope}[shift={(#1, -#2)}]
    \draw[line width=0.5pt] (-0.3, -0.3) -- (-0.3, 0.3) -- (0.3, 0.3) -- (0.3, -0.3);
    \draw[#3, fill=#3] (-0.3, -0.3) circle [radius=1pt];
    \draw[#4, fill=#4] (-0.3, 0.3) circle [radius=1pt];
    \draw[#5, fill=#5] (0.3, 0.3) circle [radius=1pt];
    \draw[#6, fill=#6] (0.3, -0.3) circle [radius=1pt];
    \end{scope}
}
\DeclareRobustCommand{\stirling}{\genfrac\{\}{0pt}{}}
\def\multichoose#1#2{\ensuremath{\left(\kern-.3em\left(\genfrac{}{}{0pt}{}{#1}{#2}\right)\kern-.3em\right)}}


\begin{document}
\section*{\raggedright 945: Plethystic formulas for ordinary, surjective and stable complex symmetric-group restriction coefficients and their inverses}
\noindent\textbf{Author:} Lee, Mitchell\par\noindent\textbf{Source:} The Frobenius transform of a symmetric function. Algebraic Combinatorics 7 (2024), publisher version, DOI 10.5802/alco.365. \url{https://alco.centre-mersenne.org/articles/10.5802/alco.365/}
\par\noindent\textbf{License:} CC BY 4.0, \url{https://creativecommons.org/licenses/by/4.0/}. The original publisher PDF cover has the explicit license notice. See the saved source/license records.
\par\noindent\textbf{Editorial material note (not author text):} The selected problem is original Theorem4.2 only, all five ordinary/surjective/inverse/stable restriction coefficient identities as one unit. Complete Sections1--4 retain finite-set category modules, image factorization, Hall adjoints, inverse plethysm and horizontal-strip stabilization with the entire human proof. Cadogan plethystic inverse and Negation Rule remain named standard prerequisites. The exact two original unsized double-bracket pairs at native202--203 use the approved bounded typography resolver. The original ytableau import is genuinely unused in these retained Sections1--4. Approved finite typography note: literal unsized double-square bracket controls resolve only through installed LatinModernMath U+27E6/U+27E7; the original stmary10 shape/width/vertical placement can differ. Every original bracket token and adjacent exponent/subscript is preserved, with all-use source/output glyph inventory. No unavailable package is acquired or emulated. All proof text and formulas below are editable author TeX. Mathematical wording, language and proof order are retained. Only the article class, title matter and bibliography presentation are adapted for independent compilation; no proof has been generated or repaired.
\par\noindent\textbf{Editorial difficulty:} H2: The proof realizes the Frobenius transform through finite-set functor categories, separates surjections by image, and transports inverse plethysm through adjoint operators before stabilizing horizontal-strip multiplicities. The complete finite-set-module construction and all five coefficient identities are preserved as one unit rather than splitting equivalent inverse formulas.
\medskip\hrule\medskip
\noindent\textbf{Author statement, complete proof and local supporting text follow in source order.}
\section{Introduction}
Let $n \geq 0$ and let~$\lambda$ be a partition with at most~$n$ parts. There is a corresponding irreducible $GL_n(\K)$-module: the Schur module $\mathbb{S}^\lambda \K^n$. Because the symmetric group $\Sym_n$ embeds in $GL_n(\K)$ by permutation matrices, one may ask: how does the restriction of $\mathbb{S}^\lambda \K^n$ to $\Sym_n$ decompose into irreducible $\Sym_n$-modules?

In other words, let~$\lambda$ and~$\mu$ be partitions and let $n = |\mu|$. What is the value of the \textit{restriction coefficient} \[r_{\lambda}^\mu = \dim \operatorname{Hom}_{\Sym_n}(V_\mu, \mathbb{S}^\lambda \K^n),\] where $V_\mu$ is the Specht module corresponding to the partition~$\mu$? This problem, called the \emph{restriction problem}, has held considerable recent interest~\cite{MR4055931, MR4356269, MR4311085, MR4295089, MR4127906}. However, there remains no known combinatorial interpretation for $r_\lambda^\mu$.

Let~$\Lambda$ be the ring of symmetric functions in the variables $x_1, x_2, x_3, \ldots$ and let $\overline{\Lambda}$ be the ring of symmetric power series in $x_1, x_2, x_3, \ldots$. In this paper, we will consider the abelian group homomorphism $\Fa \colon \Lambda \to \overline{\Lambda}$ defined on the basis $\{s_\lambda\}$ of Schur functions by \[\F{s_\lambda} = \sum_{\mu} r_{\lambda}^\mu s_\mu.\] Equivalently, $\F{s_\lambda}$ is the result of applying the Frobenius character map to the representation $\bigoplus_{n} \mathbb{S}^\lambda \K^n$ of $\bigoplus_{n} \K[\Sym_n]$. For this reason, we call~$\Fa$ the \emph{Frobenius transform}. It encodes all information about all the restriction coefficients.

\Cref{section:basics} will cover the basic properties of the Frobenius transform, many of which are implicit in the work of Orellana and Zabrocki. For example, for any symmetric functions $f, g$, we have $\F{fg} = \F{f} \ast \F{g}$, where~$\ast$ is the Kronecker product of symmetric functions. Moreover, for any $f \in \Lambda$, there exists $\FSur{f} \in \Lambda$ with the same degree and leading term as~$f$ such that $\F{f} = \FSur{f} \cdot (1 + h_1 + h_2 + \cdots)$. We refer to the map $\Fa_{\Sur} \colon \Lambda \to \Lambda$ as the \emph{surjective Frobenius transform}. Because the surjective Frobenius transform preserves degree and leading term, it has an inverse, which we denote by $\Fa_{\Sur}^{-1}$.

The \emph{induced trivial character basis} $\{\tilde{h}_\lambda\}_\lambda$ and the \emph{irreducible character basis} $\{\tilde{s}_\lambda\}_\lambda$, which were introduced by Orellana and Zabrocki~\cite{MR4295089} in 2021, can be defined in terms of the inverse surjective Frobenius transform. Namely, $\tilde{h}_\lambda = \FSurInv{h_\lambda}$ and $\tilde{s}_\lambda = \FSurInv{(1 + h_1 + h_2 + \cdots)^{\perp} s_\lambda}$, where $f^\perp\colon \Lambda \to \Lambda$ denotes the operator adjoint under the Hall inner product to multiplication by~$f$.

In \Cref{section:stable}, we will use the Frobenius transform to study \emph{stable restriction coefficients}; that is, the limits \[a_\lambda^\mu = \lim_{n \to \infty} r_{\lambda}^{(n - |\mu|, \mu_1, \ldots, \mu_{\ell(\mu)})},\] which exist for all $\lambda, \mu$ by a classical result of Littlewood~\cite{MR1576896}. In 2019, Assaf and Speyer~\cite{MR4055931} found a formula for $a_\lambda^\mu$ and for the entries $b_\lambda^\mu$ of the inverse matrix $[b_\lambda^\mu] = [a_\lambda^\mu]^{-1}$. We will provide an alternative proof of these formulas using the Frobenius transform. In \Cref{theorem:summary}, we will broadly summarize the known relationships between the five kinds of restriction coefficients considered in this paper.

In Section~5, we will prove the following theorem, which shows how to write $\Fa_{\Sur} \colon \Lambda \to \Lambda$ as a sum of operators of the form $fg^\perp$.
\begin{restatable}{theo}{expansion}\label{theorem:expansion}
    Let~$f$ be a symmetric function. Then
    \[
        \FSur{f} = \sum_{\lambda} s_\lambda (s_\lambda[h_2 + h_3 + h_4 + \cdots])^\perp f,
    \]
    where the sum is over all partitions~$\lambda$.
\end{restatable}
Since \Cref{theorem:expansion} involves plethysm and there is no known simple formula for the plethysm of Schur functions, it is not well-suited for general computation. We will, however, use it to prove that many restriction coefficients $r_\lambda^\mu$ and stable restriction coefficients $a_\lambda^\mu$ vanish.
\begin{restatable}{theo}{vanishing}\label{theorem:vanishing}
    Let $\lambda, \mu$ be partitions. If $r_\lambda^\mu > 0$, then $|\lambda \cap \hat \mu| \geq 2|\hat \mu| - |\lambda|$, where $\hat \mu = (\mu_2, \ldots, \mu_{\ell(\mu)})$ is the partition formed by removing the first part of~$\mu$.
\end{restatable}
\begin{restatable}{theo}{stablevanishing}\label{theorem:stablevanishing}
    Let $\lambda, \mu$ be partitions. If $a_\lambda^\mu > 0$, then $|\lambda \cap \mu| \geq 2|\mu| - |\lambda|$.
\end{restatable}

In Section~6, we will compute $\FSur{f}$ when~$f$ is a complete homogeneous, elementary, or power sum symmetric function. We have $\FSur{e_n} = e_n$ and $\FSur{p_n} = \sum_{d \mid n} p_d$ for $n \geq 1$. More generally, with $\N = \{0, 1, 2, \ldots\}$:

\begin{restatable}{theo}{computation}\label{theorem:computation}
Let~$\lambda$ be a partition and let $\ell = \ell(\lambda)$ be its length.
\begin{enumerate}[label=(\alph*)]
    \item \label{item:fsurh} Then \[\FSur{h_\lambda} = \sum_{M} \prod_{j \in \N^\ell} h_{M(j)},\] where the sum is over all functions $M\colon \N^\ell \to \N$ such that $M(0, \ldots, 0) = 0$ and $\sum_{j \in \N^\ell} j_i M(j) = \lambda_i$ for $i = 1, \ldots, \ell$.
    \item \label{item:fsure} Then \[\FSur{e_\lambda} = \sum_{M} \prod_{j \in \{0, 1\}^\ell}\begin{cases}
    h_{M(j)} & \mbox{if $j_1 + \cdots + j_\ell$ is even;} \\
    e_{M(j)} & \mbox{if $j_1 + \cdots + j_\ell$ is odd,}
    \end{cases}\] where the sum is over all functions $M \colon \{0, 1\}^\ell \to \N$ such that $M(0, \ldots, 0) = 0$ and $\sum_{j \in \{0, 1\}^\ell} j_i M(j) = \lambda_i$ for $i = 1, \ldots, \ell$.
    \item \label{item:fsurp} Then \[\FSur{p_\lambda} = \sum_{\pi} \prod_{U \in \pi}\left(\sum_{d \mid \gcd\{\lambda_i \colon i \in U\}} d^{|U| - 1} p_d\right),\] where the outer sum is over all partitions~$\pi$ of $\{1, \ldots, \ell\}$ into nonempty sets.
\end{enumerate}
\end{restatable}

A statement equivalent to parts (a) and (b) of this theorem has appeared previously in the work of Orellana and Zabrocki \cite[Equation~(6)]{MR4245137}.

\Cref{theorem:computation} has the following interesting consequence. For any partition~$\mu$, denote by $D(\mu)$ the size of the Durfee square of~$\mu$; that is, $D(\mu)$ is the largest integer~$d$ such that $\mu_{d} \geq d$ \cite[Chapter~8]{MR2122332}.

\begin{restatable}{theo}{durfeebound}\label{theorem:durfeebound}
    Let~$\mu$ be a partition and let $k \geq 1$ be an integer. The following are equivalent:
    \begin{enumerate}[label=(\Alph*)]
        \item There exists a partition~$\lambda$ such that $\lambda_1 \leq k$ and $r_{\lambda}^\mu > 0$.
        \item $D(\mu) \leq 2^{k - 1}$.
    \end{enumerate}
\end{restatable}

In particular, $r_{\lambda}^\mu = 0$ if $D(\mu) > 2^{\lambda_1 - 1}$.

Finally, in Section~7, we will compute $\FSurInv{e_\lambda}$ and $\FSurInv{h_\lambda}$ (Theorem~7.7). In particular, we will prove that \[\FSurInv{e_\lambda} = \sum_{\mu} (-1)^{|\lambda| - |\mu|} L_\lambda^\mu e_\mu,\] where $L_\lambda^\mu$ is a nonnegative integer with an explicit combinatorial interpretation involving Lyndon words (Corollary~7.9).

\section{Preliminaries}\label{section:preliminaries}
Apart from the definition of the restriction coefficients $a_\lambda^\mu$ (\Cref{definition:restriction} below), all the definitions in this section can be found in any standard reference on the theory of symmetric functions such as~\cite{MR1464693},~\cite[Chapter~I]{MR3443860}, or~\cite[Chapter~7]{MR4621625}.

Let $\Lambda = \Lambda_\mathbb{Z}$ denote the ring of symmetric functions over $\mathbb{Z}$ in the variables $x_1, x_2, x_3, \ldots$. For $n \geq 0$, let $\Lambda_n$ denote the subgroup of~$\Lambda$ consisting of all symmetric functions that are homogeneous of degree~$n$. Let $\overline{\Lambda}$ denote the ring of symmetric formal power series (that is, formal sums $\sum_{n} f_n$, where each $f_n \in \Lambda_n$). Let~$\langle \cdot, \cdot \rangle \colon \Lambda \times \overline{\Lambda} \to \Z$ denote the Hall inner product.

For any partition $\lambda = (\lambda_1, \ldots, \lambda_\ell)$, define the \emph{length} $\ell(\lambda) = \ell$ and the \emph{size} $|\lambda| = \lambda_1 + \cdots + \lambda_\ell$. Let $m_i(\lambda)$ be the number of times the part~$i$ appears in~$\lambda$, and define $z_\lambda = \prod_i i^{m_i(\lambda)} (m_i(\lambda))!$. Let $\lambda^T$ denote the dual (\ie transpose) of~$\lambda$. Let $m_\lambda, e_\lambda, h_\lambda, p_\lambda, s_\lambda \in \Lambda$ denote the monomial, elementary, homogeneous, power sum, and Schur symmetric functions respectively. Let~$V_\lambda$ denote the corresponding \emph{Specht module}, which is an irreducible $\Sym_{|\lambda|}$-module. Let~$\chi_\lambda$ denote the character of $V_\lambda$. Let~$\mathbb{S}^\lambda$ denote the corresponding \emph{Schur functor}, which is an endofunctor of the category of vector spaces over~$\K$.

For any partitions $\lambda, \mu$, define the \emph{intersection} $\lambda \cap \mu$ by $\ell(\lambda \cap \mu) = \min(\ell(\lambda), \ell(\mu))$ and $(\lambda \cap \mu)_i = \min(\lambda_i, \mu_i)$. That is, it is the partition whose Young diagram is the intersection of the Young diagrams of~$\lambda$ and~$\mu$. Let us say that $\lambda \subseteq \mu$ if $\lambda \cap \mu = \lambda$.

Let $\omega \colon \Lambda \to \Lambda$ be the ring homomorphism given by $\omega(p_k) = (-1)^{k-1} p_k$. Recall that~$\omega$ is an involution and for all partitions~$\lambda$, we have $\omega(h_\lambda) = e_\lambda$ and $\omega(s_\lambda) = s_{\lambda^T}$.

The \emph{Lyndon symmetric function} $L_n$ is given by \[L_n = \frac{1}{n} \sum_{d \mid n} \mu(d) p_d^{n/d}\] for $n \geq 1$, where $\mu \colon \{1, 2, 3, \ldots\} \to \{-1, 0, 1\}$ is the M\"obius function.

For all $f \in \overline{\Lambda}$, the \emph{skewing operator} $f^\perp \colon \Lambda \to \Lambda$ is the adjoint to multiplication by~$f$ under the Hall inner product: \[\langle g, f^\perp h \rangle = \langle fg, h \rangle.\]

We say that a Schur function $s_\lambda$ \emph{appears} in $f \in \overline{\Lambda}$ if $\langle s_\lambda, f \rangle \neq 0$. We say that~$f$ is \emph{Schur positive} if $\langle s_\lambda, f \rangle \geq 0$ for all partitions~$\lambda$.

Let $C_n$ denote the space of all~$\K$-valued class functions on $\Sym_n$. Let $R_n$ denote the additive group of all virtual characters on $\Sym_n$. In other words, $R_n$ is the subgroup of~$C_n$ generated by the irreducible characters $\chi_\lambda$. The~$n$th \emph{Frobenius character map} is the map $\ch_n \colon R_n \to \Lambda_n$ defined by \[\ch_n(\chi) = \frac{1}{n!}\sum_{w \in \Sym_n} \chi(w) p_{c(w)},\] where $c(w)$ denotes the cycle type of~$w$. It is well-known that $\ch_n$ is an isomorphism and that $\ch_n(\chi_\lambda) = s_\lambda$ for all~$\lambda$ with $|\lambda| = n$.

The \emph{Kronecker product} is the unique bilinear operator $\ast \colon \Lambda \times \Lambda \to \Lambda$ satisfying $p_\lambda \ast p_\mu = \delta_{\lambda \mu} z_\lambda p_\lambda$ for all $\lambda, \mu$. It extends to a bilinear operator $\ast \colon \overline{\Lambda} \times \overline{\Lambda} \to \overline{\Lambda}$ by continuity. The Frobenius character map and the Kronecker product are related in the following way: for any $\chi_1, \chi_2 \in R_n$, we have $\ch(\chi_1 \chi_2) = \ch(\chi_1) \ast \ch(\chi_2)$.

For any $f, g \in \overline{\Lambda}$, let $f[g]$ denote the plethysm of~$f$ by~$g$. This is well-defined as long as $f \in \Lambda$ or~$g$ has no constant term.

\begin{prop}[{Plethystic Addition Formula, \cite[Section~3.2]{MR2765321}}]\label{proposition:plethysticaddition}
    Let~$\lambda$ be a partition and let $f, g \in \overline{\Lambda}$. Then
    \[s_{\lambda}[f + g] = \sum_{\mu} s_{\lambda / \mu}[f] s_\mu[g],\] where the sum is over all partitions~$\mu$.
\end{prop}

For a variable $t$, let
\begin{alignat*}{4}
H(t) &= \sum_{n \geq 0} h_n t^n &&= \prod_{i \geq 1} \frac{1}{1 - x_i t} &&= \exp\left(\sum_{k \geq 1} \frac{p_k}{k}t^k\right)&&\in \Lambda \llbracket t \rrbracket \\
E(t) &= \sum_{n \geq 0} e_n t^n &&= \prod_{i \geq 1} (1 + x_i t) &&= \exp\left(\sum_{k \geq 1} \frac{p_k}{k} (-1)^{k-1} t^k\right) &&\in \Lambda \llbracket t \rrbracket.
\end{alignat*}
It is clear that $E(t) = \frac{1}{H(-t)}$. Let $H = H(1) = 1 + h_1 + h_2 + \cdots$ and $H_+ = h_1 + h_2 + h_3 + \cdots = H - 1$.

For any partitions $\lambda, \mu, \nu$, the \emph{Littlewood--Richardson coefficient} $c^\nu_{\lambda\mu}$ is the Hall inner product $\langle s_\nu, s_\lambda s_\mu \rangle$. Define $g_{\lambda\mu\nu} = \langle s_\nu, s_\lambda \ast s_\mu \rangle$.

\begin{defi}\label{definition:restriction}
    Let $\lambda, \mu$ be partitions and let $n = |\mu|$. Then, the Schur module $\mathbb{S}^\lambda \K^n$ can be considered as an $\Sym_n$-module, with $\Sym_n$ acting on $\K^n$ by permutation matrices. The \emph{restriction coefficient} \[r_\lambda^\mu = \dim \operatorname{Hom}_{\Sym_n}(V_\mu, \mathbb{S}^\lambda \K^n)\] is the multiplicity of the Specht module $V_\mu$ in $\mathbb{S}^\lambda \K^n$. 
\end{defi}


\section{The Frobenius transform: definition and basic properties}
\label{section:basics}
Recall from the introduction that it is a long-standing open problem to find a combinatorial interpretation of $r_\lambda^\mu$. As a potential way to approach this problem, we now define the Frobenius transform, which is the primary object of study in this paper.
\begin{defi}
    The \emph{Frobenius transform} is the abelian group homomorphism $\Fa \colon \Lambda \to \overline{\Lambda}$ defined on the basis $\{s_\lambda\}$ by \[\F{s_\lambda} = \sum_\mu r_{\lambda}^\mu s_\mu,\] where the sum is over all partitions~$\mu$.
\end{defi}
\begin{rema}\label{remark:halladjoint}
By a classical result of Littlewood~\cite{MR1576896}, we have \[r_\lambda^\mu = \langle s_\lambda, s_\mu[H] \rangle.\] Hence,~$\Fa$ is adjoint to plethysm by~$H$ under the Hall inner product.
\end{rema}
Here is the reason for calling~$\Fa$ the Frobenius transform. Let $n \geq 0$ and let~$\lambda$ be any partition. Then $\mathbb{S}^\lambda \K^n$, considered as an $\Sym_n$-module, can be expressed as a direct sum of Specht modules:
\[\bigoplus_{|\mu| = n} r_{\lambda}^\mu V_\mu = \mathbb{S}^\lambda \K^n.\]
Taking the character of both sides and applying the Frobenius character map, we obtain \begin{equation}\label{equation:fslambdan1}\sum_{|\mu| = n} r_{\lambda}^\mu s_\mu = \ch_n(\chi_{\mathbb{S}^\lambda \K^n}).\end{equation} In other words, the degree~$n$ part of $\F{s_\lambda}$ is equal to the Frobenius character of~$\mathbb{S}^\lambda \K^n$.

\begin{exam}\label{example:fer}
    Let $r \geq 0$. We will compute $\F{e_r}$. First, $e_r = s_{\lambda}$, where $\lambda = (1^r)$. Hence, for any~$n$, the degree~$n$ part of $\F{e_r}$ is the Frobenius character of \[\mathbb{S}^{\lambda} \K^n = \wedge^r \K^n = \operatorname{Ind}_{\Sym_r \times \Sym_{n - r}}^{\Sym_n} (V_{(1^r)} \otimes V_{(n-r)}),\] considered as an $\Sym_n$-module. Thus, it is equal to $e_r h_{n - r}$ \cite[Proposition~7.18.2]{MR4621625}. Taking the sum over all~$n$ yields $\F{e_r} = e_r \cdot H$.
\end{exam}

\subsection{The Frobenius transform and representations of combinatorial categories}\label{subsection:category}
The purpose of this subsection is to provide an alternate perspective on the Frobenius transform. This subsection is not essential to the proofs of our main results, so the reader may skip it. For a more complete introduction to the representation theory of categories, see Wiltshire-Gordon's 2016 PhD thesis~\cite{MR3641124}.

In what follows, let $\Vect_\K$ be the category whose objects are vector spaces over~$\K$ and whose morphisms are linear transformations. (The ground field~$\K$ can be replaced by any algebraically closed field of characteristic $0$.)

\begin{defi}
    Let $\mathcal{C}$ be a category. A \emph{$\mathcal{C}$-module} (over the ground field~$\K$) is a functor $M[\bullet] \colon \mathcal{C} \to \Vect_\K$. The \emph{category of $\mathcal{C}$-modules} is the functor category $\Mod^\mathcal{C} = (\Vect_\K)^\mathcal{C}$.
    
    We say that a $\mathcal{C}$-module~$M$ is \emph{finite-dimensional} if $M[U]$ is finite-dimensional for all $U \in \operatorname{Ob}(\mathcal{C})$.

    If~$M$ and~$N$ are $\mathcal{C}$-modules, we may form the \emph{direct sum} $M \oplus N$ by $(M \oplus N)[U] = M[U] \oplus N[U]$.
\end{defi}

Let $\rmBij$ be the category whose objects are finite sets and whose morphisms are bijections, and let~$M$ be a $\rmBij$-module. (Previous authors have referred to $\rmBij$-modules as \emph{linear species} or \emph{tensor species}~\cite{MR0927763, MR1132426}.) For $n \geq 0$, denote by $M[n]$ the vector space $M[[n]] = M[\{1, \ldots, n\}]$, which is an $\Sym_n$-module. Then~$M$ is uniquely determined, up to natural isomorphism, by the sequence $M[0], M[1], M[2], \ldots$ of symmetric group modules; namely, 
\begin{equation}\label{equation:bmodule}
M[U] = \K \Bij([n], U) \otimes_{\Sym_n} M[n]
\end{equation} for all finite sets~$U$ with $|U| = n$.
Additionally, if~$M$ is finite-dimensional, then we may form the series
\[\ch(M) = \sum_{n} \ch_n(\chi_{M[n]}) \in \overline{\Lambda},\] which we call the \emph{Frobenius character of~$M$}. It also uniquely determines~$M$ up to natural isomorphism.

Not every symmetric power series can be written in the form $\ch(M)$, where~$M$ is a finite-dimensional $\rmBij$-module. (Every symmetric power series can be written as the Frobenius character of a \emph{virtual $\rmBij$-module}, but we will not define virtual $\rmBij$-modules in this article.) However, it is still helpful to think of~$\ch$ as a partial correspondence between (isomorphism classes of) finite-dimensional $\rmBij$-modules and symmetric power series. Many concepts from the theory of symmetric power series have analogous concepts in the theory of $\rmBij$-modules. For example, in \Cref{definition:speciesproduct} below, we will define the \emph{product} $M \cdot N$ of two $\rmBij$-modules $M, N$. Via the Frobenius character, this is analogous to the product of symmetric power series in the sense that $\ch(M \cdot N) = \ch(M) \ch(N)$ (\Cref{proposition:moduleproduct}).

In what follows, we will find the construction in the theory of $\rmBij$-modules which is analogous to the Frobenius transform. More precisely, let~$M$ be a $\rmBij$-module such that $M_n = 0$ for all but finitely many~$n$. Then $\ch(M)$ is in fact a symmetric function, so its Frobenius transform $\F{\ch(M)}$ is well-defined. We will construct a $\rmBij$-module whose Frobenius character is $\F{\ch(M)}$. Before we do, we need the following definitions.

\begin{defi}[{\cite[Definition~2.6.1]{MR3641124}}]
    Let $F \colon \mathcal{C} \to \mathcal{D}$ be a functor. The \emph{pullback functor} $F^* \colon \Mod^\mathcal{D} \to \Mod^\mathcal{C}$ is given by $F^*M = M \circ F$.
\end{defi}
\begin{defi}[{\cite[Proposition~2.6.2]{MR3641124}}]
    Let $F \colon \mathcal{C} \to \mathcal{D}$ be a functor. The \emph{left Kan extension functor} $F_! \colon \Mod^\mathcal{C} \to \Mod^\mathcal{D}$ is the left adjoint to the pullback $F^*$, if it exists:
    \[\Hom_{\Mod^\mathcal{D}}(F_! M, N) = \Hom_{\Mod^\mathcal{C}}(M, F^*N).\]
\end{defi}

Let $\rmFun$ be the category whose objects are finite sets and whose morphisms are functions. Clearly, $\rmBij$ is a subcategory of $\rmFun$; let $\iota \colon \rmBij \to \rmFun$ be the inclusion functor.

\begin{prop}\label{proposition:leftkanextension}
    The left Kan extension functor $\iota_! \colon \Mod^\rmBij \to \Mod^\rmFun$ exists. Moreover, for every $\rmBij$-module~$M$, the left Kan extension $\iota_! M$ is given on finite sets~$U$ by \[(\iota_! M)[U] = \bigoplus_{n} \K U^n \otimes_{\Sym_n} M[n].\] Here $\K U^n \otimes_{\Sym_n} M[n]$ denotes the quotient of the tensor product $\K U^n \otimes M[n]$ by the relation that $a \otimes w b \sim aw \otimes b$ for all $a \in \K U^n$, $b \in M_n$, and $w \in \Sym_n$.
\end{prop}
\begin{proof}
    This follows from \cite[Proposition~2.6.7]{MR3641124}.
\end{proof}
\begin{rema}
    Joyal refers to $\iota_! M \colon \rmFun \to \Vect_\K$ as the \emph{analytic functor} corresponding to the tensor species~$M$ \cite[Definition~4.2]{MR0927763}.
\end{rema}
We are now ready to show that $\iota^* \iota_!$ is the construction analogous to the Frobenius transform. 
\begin{prop}\label{proposition:kanfrobenius}
    Let~$M$ be a finite-dimensional $\rmBij$-module such that $M[n] = 0$ for all but finitely many~$n$. Then $\iota^*\iota_! M$ is finite-dimensional and $\F{\ch(M)} = \ch(\iota^*\iota_! M)$.
\end{prop}
\begin{proof}
    For any partition~$\lambda$, define the $\rmBij$-module $M_\lambda$ by \[(M_\lambda)[n] = \begin{cases}
        V_\lambda & \mbox{if $n = |\lambda|$;}\\
        0 & \mbox{otherwise,}
    \end{cases}\] extending to all of $\rmBij$ using \eqref{equation:bmodule}.
    
    We have that~$M$ can be written as a direct sum of the $M_\lambda$. Since $\iota^*$ and $\iota_!$ preserve direct sums, it is enough to prove the proposition for $M = M_\lambda$. In this case, $\ch(M) = \ch_n(\chi_\lambda) = s_\lambda$.

    On the other hand, by \Cref{proposition:leftkanextension}, we have for $m \geq 0$ that 
    \begin{align*}
        (\iota^*\iota_! M)[m] &= \bigoplus_{n} \K [m]^n \otimes_{\Sym_n} M[n] \\
        &= \bigoplus_{n} (\K^m)^{\otimes n} \otimes_{\Sym_n} M[n] \\
        &= (\K^m)^{\otimes |\lambda|} \otimes_{\Sym_{|\lambda|}} V_\lambda.
    \end{align*}
    By Schur-Weyl duality, this is isomorphic as a $GL_m$-module to $\mathbb{S}^\lambda \K^m$. Hence, as a $\Sym_m$-module, it decomposes into irreducibles as
    \[(\iota^* \iota_! M)[m] = \bigoplus_{|\mu| = m} r^\mu_\lambda V_\mu.\] It follows that \[\ch(\iota^* \iota_! M) = \sum_{\mu} r^{\mu}_\lambda s_\mu = \F{s_\lambda} = \F{\ch(M)}\] as desired.
\end{proof}
We will occasionally make use of the following definition in later remarks.
\begin{defi}[{\cite[Section~4.1]{MR0927763}}]\label{definition:speciesproduct}
    Let $M, N$ be $\rmBij$-modules. Define the \emph{product} $M \cdot N$ to be the $\rmBij$-module given by \[(M \cdot N)[U] = \bigoplus_{\substack{U_1 \cup U_2 = U \\ U_1 \cap U_2 = \emptyset}} M[U_1] \otimes N[U_2].\]
\end{defi}
\begin{prop}[{\cite[Proposition~2.1]{MR1132426}}]
    \label{proposition:moduleproduct}
    Let~$M$,~$N$ be finite-dimensional $\rmBij$-modules. Then \[\ch(M \cdot N) = \ch(M) \ch(N).\]
\end{prop}
\subsection{The Frobenius transform and evaluation at roots of unity}\label{subsection:rootsofunity}
Let us now describe the expansion of $\F{s_\lambda}$ in the power sum basis. In order to simplify the description, we will do it one degree at a time. In this subsection, we will write $f_n$ to mean the degree~$n$ part of~$f$.

By \eqref{equation:fslambdan1}, we have
\begin{equation}\label{equation:fslambdan2}
(\F{s_\lambda})_n = \ch_n(\chi_{\mathbb{S}^\lambda \K^n}) = \frac{1}{n!} \sum_{w \in \Sym_n} \chi_{\mathbb{S}^\lambda \K^n}(w) p_{c(w)}.
\end{equation}
It is well-known \cite[Section~8.3]{MR1464693} that the character of $\mathbb{S}^\lambda \K^n$ as a $GL_n$-module is the Schur function $s_\lambda$ itself. In other words, if $g \in GL_n(\K)$ has eigenvalues $x_1, \ldots, x_n$, then $\chi_{\mathbb{S}^\lambda \K^n}(g)$ is equal to the evaluation $s_\lambda(x_1, \ldots, x_n)$.

For $w \in \Sym_n$, we have $\chi_{\mathbb{S}^\lambda \K^n}(w) = s_\lambda(x_1, \ldots, x_n)$, where $x_1, \ldots, x_n$ are the eigenvalues of the permutation matrix $P_w$. Let $\mu = (\mu_1, \ldots, \mu_\ell)$ be the cycle type of~$w$. Then the eigenvalues of $P_w$ are the roots of unity
\begin{gather*}
    1, \exp\left(\frac{2 \pi i}{\mu_1}\right), \exp\left(\frac{4 \pi i}{\mu_1}\right), \ldots, \exp\left(\frac{2(\mu_1 - 1) \pi i}{\mu_1}\right), \\
    \vdots \\
    1, \exp\left(\frac{2 \pi i}{\mu_\ell}\right), \exp\left(\frac{4 \pi i}{\mu_\ell}\right), \ldots, \exp\left(\frac{2(\mu_\ell - 1) \pi i}{\mu_\ell}\right).
\end{gather*}
Following Orellana and Zabrocki~\cite{MR4275829, MR4295089}, let $\Xi_\mu \in \K^n$ denote this sequence. We have that $\chi_{\mathbb{S}^\lambda \K^n}(w)$ is the result of evaluating $s_\lambda$ at these roots of unity: \[\chi_{\mathbb{S}^\lambda \K^n}(w) =
s_\lambda\left(\Xi_\mu\right).\]
Now, let us group the terms on the right-hand side of \eqref{equation:fslambdan2} according to the cycle type $\mu = c(w)$. The number of permutations $w \in \Sym_n$ with cycle type~$\mu$ is $\frac{n!}{z_\mu}$, so 
\begin{align*}(\F{s_\lambda})_n &= \frac{1}{n!} \sum_{w \in \Sym_n} \chi_{\mathbb{S}^\lambda \K^n}(w) p_{c(w)} \\
&= \frac{1}{n!} \sum_{|\mu| = n} \frac{n!}{z_\mu}s_\lambda(\Xi_\mu)p_\mu \\
&= \sum_{|\mu| = n} s_\lambda(\Xi_\mu) \frac{p_\mu}{z_\mu}.\end{align*}
Taking the sum over all~$n$, we obtain \[\F{s_\lambda} = \sum_{\mu} s_\lambda(\Xi_\mu) \frac{p_\mu}{z_\mu}.\]
Finally, by linearity, we may extend this result to any $f \in \Lambda$. We have proved the following.
\begin{prop}\label{proposition:rootsofunityformula}
    Let $f \in \Lambda$. Then \[\F{f} = \sum_{\mu} f(\Xi_\mu) \frac{p_\mu}{z_\mu}.\]
\end{prop}
In the notation of Orellana and Zabrocki~\cite{MR4275829}, this proposition can be written as \[\F{f} = \phi_0(f) + \phi_1(f) + \phi_2(f) + \cdots.\]

\subsection{The Frobenius transform and the Kronecker product}
The Frobenius transform relates the ordinary product of symmetric functions to the Kronecker product in the following way.
\begin{prop}[{cf. \cite[Section~2.3]{MR4275829}}]
\label{proposition:kronecker}
Let $f, g \in \Lambda$. Then $\F{fg} = \F{f} \ast \F{g}$.
\end{prop}
We provide two different proofs of \Cref{proposition:kronecker}: a category-theoretic proof using \Cref{proposition:kanfrobenius} and a direct computational proof using \Cref{proposition:rootsofunityformula}.
\begin{proof}[First proof]
    Because both sides of the desired equation are bilinear, we may assume that there exist $\rmBij$-modules $M, N$ such that $\ch(M) = f$ and $\ch(N) = g$. By \Cref{proposition:moduleproduct}, we have $\ch(M \cdot N) = \ch(M) \ch(N) = fg$, where $M \cdot N$ is the product as defined in \Cref{definition:speciesproduct}.
    
    By \cite[Equation~2.1(ii)]{MR0927763}, we have $\iota^* \iota_! (M \cdot N) = (\iota^* \iota_! M) \otimes_{\K} (\iota^* \iota_! N)$, where $\otimes_\K$ denotes the object-wise tensor product of $\rmBij$-modules. Applying~$\ch$ to both sides, we obtain $\F{fg} = \F{f} \ast \F{g}$, as desired.
\end{proof}
\begin{proof}[Second proof]
By \Cref{proposition:rootsofunityformula}, we have
\begin{align*}
    \F{f} \ast \F{g} &= \left(\sum_{\mu} f(\Xi_\mu) \frac{p_\mu}{z_\mu}\right) \ast \left(\sum_{\mu} g(\Xi_\mu) \frac{p_\mu}{z_\mu}\right) \\
    &= \sum_{\mu} f(\Xi_\mu)g(\Xi_\mu) \frac{p_\mu}{z_\mu} \\
    &= \F{fg}
\end{align*}
as desired.
\end{proof}

\begin{coro}
\label{corollary:productcoefficients}
Let $\lambda, \mu, \nu$ be partitions. Then
\[\sum_{\nu'} r^{\nu}_{\nu'} c^{\nu'}_{\lambda\mu} = \sum_{\lambda', \mu'} r^{\lambda'}_{\lambda} r^{\mu'}_{\mu} g_{\lambda'\mu'\nu}.\]
\end{coro}
\begin{proof}
In \Cref{proposition:kronecker}, take $f = s_\lambda$ and $g = s_\mu$. Then take the Hall inner product of both sides with $s_\nu$.
\end{proof}

\subsection{The surjective Frobenius transform}
\begin{prop}
\label{proposition:fsur}
Let $f \in \Lambda$. Then there exists a symmetric function $\FSur{f} \in \Lambda$ such that $\F{f} = \FSur{f} \cdot H$. Moreover, $\FSur{f}$ has the same degree and leading term as~$f$.
\end{prop}
For example, in \Cref{example:fer} we showed that $\F{e_r} = e_r \cdot H$. Thus, $\FSur{e_r} = e_r$. (Note, however, that in general, $\Fa_{\Sur}$ does not preserve the property of being homogeneous.) Again, we provide two separate proofs of this proposition: a category-theoretic proof and a direct computational proof.
\begin{proof}[First proof]
Because both sides of the desired equation are linear in~$f$, we may assume that there exists a $\rmBij$-module~$M$ such that $\ch(M) = f$.

Let $\rmSur$ be the category whose objects are finite sets and whose morphisms are surjections, and let $\kappa \colon \rmBij \to \rmSur$ be the inclusion functor. Let $\kappa_! \colon \Mod^\rmBij \to \Mod^\rmFun$ be the left Kan extension functor along~$\kappa$, which exists by \cite[Proposition~2.6.7]{MR3641124}. We claim that the choice $\FSur{f} = \ch(\kappa^* \kappa_!M)$ satisfies the desired properties. For this, we will show that there is a natural isomorphism
\begin{equation}\label{eq:iotakappa}
    \iota^* \iota_! M = \kappa^* \kappa_! M \cdot \K E,
\end{equation}
where~$\iota$ and~$\cdot$ are defined as in \Cref{subsection:category} and $\K E$ is the $\rmBij$-module given by $(\K E)[U] = \K$ for all finite sets~$U$.

By \Cref{proposition:leftkanextension}, we have
\begin{equation}\label{eq:iotaiota}
    (\iota^* \iota_! M)[U] = \bigoplus_n \K U^n \otimes_{\Sym_n} M[n].
\end{equation}
Think of $U^n$ as the set of all functions $[n] \to U$. By grouping those functions by their image~$V$, we may write $U^n$ as a disjoint union:
\[U^n = \sum_{V \subseteq U} \Sur([n], V).\]
Substituting into \eqref{eq:iotaiota}, we obtain
\begin{align*}
    (\iota^* \iota_! M)[U] &= \bigoplus_n \K \left(\sum_{V \subseteq U} \Sur([n], V)\right) \otimes_{\Sym_n} M[n] \\
    &= \bigoplus_n \left(\bigoplus_{V \subseteq U} \K\Sur([n], V)\right) \otimes_{\Sym_n} M[n] \\
    &= \bigoplus_{V \subseteq U} \bigoplus_n \K\Sur([n], V) \otimes_{\Sym_n} M[n] \\
    &= \bigoplus_{V \subseteq U} (\kappa^*\kappa_!M)[V]
\end{align*}
where the last step follows from \cite[Proposition~2.6.7]{MR3641124}. Recognizing the latter as $(\kappa^* \kappa_! M \cdot \K E)[U]$, we have shown \eqref{eq:iotakappa}. Applying~$\ch$ to both sides and using \Cref{proposition:kanfrobenius} and \Cref{proposition:moduleproduct} yields $\F{f} = \ch(\kappa^* \kappa_!M) \cdot H$.

It remains to show that $\ch(\kappa^* \kappa_!M)$ has the same degree and leading term as~$f$. For this, we will prove that $(\kappa^* \kappa_!M)[m]$ and $M[m]$ are isomorphic as $\Sym_m$-modules for all $m \geq \deg f$. Consider the natural isomorphism \[(\kappa^*\kappa_!M)[U] = \bigoplus_n \K\Sur([n], U) \otimes_{\Sym_n} M[n]\] with $U = [m]$. In each summand, the factor $\K\Sur([n], [m])$ vanishes when $n < m$ and the factor $M[n]$ vanishes when $n > \deg f$. If $m \geq \deg f$, this means that every term vanishes except possibly the term $n = m$. Hence,
\begin{align*}
    (\kappa^*\kappa_!M)[m] &= \K\Sur([m], [m]) \otimes_{\Sym_m} M[m] \\
    &= \K\Sym_m \otimes_{\Sym_m} M[m] \\
    &= M[m]
\end{align*}
as desired.
\end{proof}
\begin{proof}[Second proof]
We claim that the choice
\[\FSur{f} = \sum_{\mu} \langle f, s_\mu[H_+]\rangle s_\mu\] satisfies the desired properties. In other words, we may take $\Fa_{\Sur}$ to be the adjoint to plethysm by $H_+$ under the Hall inner product.

By \Cref{remark:halladjoint}, we have
\begin{align*}
    \F{f} &= \sum_{\lambda} \langle f, s_\lambda[H]\rangle s_\lambda \\
    &= \sum_{\lambda} \langle f, s_\lambda[H_+ + 1]\rangle s_\lambda.
\end{align*}
By the plethystic addition formula (\Cref{proposition:plethysticaddition}), we have
\begin{align*}
    \F{f} &= \sum_{\lambda, \mu} \langle f, s_\mu[H_+] s_{\lambda / \mu}[1]\rangle s_\lambda.
\end{align*}
We have \[s_{\lambda / \mu}[1] = s_{\lambda / \mu}(1, 0, 0, \ldots) = \begin{cases} 1 & \mbox{if $\lambda / \mu$ is a horizontal strip;} \\ 0 & \mbox{otherwise,}\end{cases}\] so 
\begin{align*}
    \F{f} &= \sum_{\substack{\lambda, \mu \\ \text{$\lambda/\mu$ h. strip}}} \langle f, s_\mu[H_+]\rangle s_\lambda \\
    &= \sum_{\mu}\langle f, s_\mu[H_+] \rangle \left(\sum_{\substack{\lambda \\ \text{$\lambda/\mu$ h. strip}}}s_\lambda\right) \\
    &= \sum_{\mu}\langle f, s_\mu[H_+] \rangle s_\mu \cdot H,
\end{align*}
where the last equality follows from the Pieri rule.

Now, the only thing left to show is that~$f$ has the same degree and leading term as \[\sum_{\mu} \langle f, s_\mu[H_+]\rangle s_\mu.\] This follows directly from the observation that if $|\mu| \geq \deg f$, then $\langle f, s_\mu[H_+]\rangle = \langle f, s_\mu\rangle$.
\end{proof}
We refer to $\Fa_{\Sur} \colon \Lambda \to \Lambda$ as the \emph{surjective Frobenius transform}. Clearly, it is invertible:
\begin{coro}
    There exists a two-sided inverse $\Fa_{\Sur}^{-1} \colon \Lambda \to \Lambda$ of $\Fa_{\Sur}$.
\end{coro}
\begin{proof}
    Define $\mathcal{M} \colon \Lambda \to \Lambda$ by $\mathcal{M}\{f\} = f - \FSur{f}$. By \Cref{proposition:fsur}, we have \[\deg(\mathcal{M}\{f\}) < \deg f\] for any $f \in \Lambda \setminus \{0\}$. Hence, $\mathcal{M}^k\{f\} = 0$ for any $k > \deg f$.

    Define $\Fa_{\Sur}^{-1} \colon \Lambda \to \Lambda$ by \[\FSurInv{f} = f + \mathcal{M}\{f\} + \mathcal{M}^2\{f\} + \cdots.\] This is well-defined by the above, and it is easy to check that it is the two-sided inverse of $\Fa_{\Sur}$.
\end{proof}

Like the ordinary Frobenius transform, the surjective Frobenius transform can be described in terms of its matrix entries in the Schur basis.
\begin{defi}
Let $\lambda, \mu$ be partitions. Define the \emph{surjective restriction coefficient} \[t_\lambda^\mu = \langle \FSur{s_\lambda}, s_\mu\rangle\] and define the \emph{inverse surjective restriction coefficient} \[u_\lambda^\mu = \langle \FSurInv{s_\lambda}, s_\mu\rangle.\]
\end{defi}
By the above, we have $t_\lambda^\mu = u_\lambda^\mu = \delta_{\lambda \mu}$ for $|\mu| \geq |\lambda|$.
\section{Stable restriction}
\label{section:stable}
Define the \emph{stable restriction coefficients} $a_\lambda^\mu$ as follows. For any partition $\mu = (\mu_1, \ldots, \mu_\ell)$ and any $n \geq \mu_1 + |\mu|$, define $\mu^{(n)} = (n - |\mu|, \mu_1, \ldots, \mu_\ell)$. Then, for any partitions $\lambda, \mu$, the stable restriction coefficient $a_\lambda^\mu$ is defined by the following limit, which exists by a classical result of Littlewood~\cite{MR1576896}: \[a_\lambda^\mu = \lim_{n \to \infty} r_{\lambda}^{\mu^{(n)}}.\] Littlewood also showed that $a_{\lambda}^\mu = \delta_{\lambda \mu}$ if $|\lambda| \leq |\mu|$, so the infinite matrix $[a_\lambda^\mu]$, with rows and columns indexed by partitions in increasing order of size, is upper unitriangular. In 2019, Assaf and Speyer~\cite{MR4055931} found the following formula for the entries $b_\lambda^\mu$ of the inverse matrix $[b_\lambda^\mu] = [a_\lambda^\mu]^{-1}$:
\begin{restatable}[{\cite[Theorem~2]{MR4055931}}]{theorem}{assafspeyer}
\label{theorem:assafspeyer}
Let $\lambda, \mu$ be partitions. Then \[b_\lambda^\mu = (-1)^{|\lambda| - |\mu|} \langle s_{\lambda^T}, s_{\mu^T}[L_1 + L_2 + L_3 + \cdots] \cdot H \rangle.\] In particular, $(-1)^{|\lambda| - |\mu|}b_\lambda^\mu$ is a nonnegative integer.
\end{restatable}
In this section, we will provide an alternative proof of \Cref{theorem:assafspeyer}. In fact, we will prove similar plethystic formulas for all five kinds of restriction coefficients defined so far. These formulas have all been collected into \Cref{theorem:summary} below.

\begin{theorem}\label{theorem:summary}
    Let $\lambda, \mu$ be partitions with $|\lambda| = m$ and $|\mu| = n$.
    \begin{enumerate}[label=(\alph*)]
        \item The restriction coefficient $r_\lambda^\mu$ is given by 
        \begin{align*}
            r_{\lambda}^\mu
            &= \langle \F{s_\lambda}, s_\mu\rangle \\
            &= \langle s_\lambda, s_\mu[H] \rangle.
        \end{align*}
        \item The surjective restriction coefficient $t_\lambda^\mu$ is given by 
        \begin{align*}
            t_{\lambda}^\mu &= \langle \FSur{s_\lambda}, s_\mu\rangle\\
            &= \langle s_\lambda, s_\mu[H_+] \rangle.
        \end{align*}
        \item The inverse surjective restriction coefficient $u_\lambda^\mu$ is given by
        \begin{align*}
            u_\lambda^\mu &= \langle \FSurInv{s_\lambda}, s_\mu \rangle \\
            &= \langle s_\lambda, s_\mu[\omega(L_1) - \omega(L_2) + \omega(L_3) - \cdots] \rangle \\
            &= (-1)^{m - n}\langle s_{\lambda^T}, s_{\mu^T}[L_1 + L_2 + L_3 + \cdots] \rangle.
        \end{align*}
        \item The stable restriction coefficient $a_\lambda^\mu$ is given by
        \begin{align*}
            a_\lambda^\mu &= \langle H^\perp \FSur{s_\lambda}, s_\mu\rangle \\
            &= \langle \FSur{s_\lambda}, s_\mu \cdot H\rangle \\
            &= \langle s_\lambda, (s_\mu \cdot H)[H_+]\rangle.
        \end{align*}
        \item The inverse stable restriction coefficient $b_\lambda^\mu$ is given by
        \begin{align*}
            b_\lambda^\mu &= \langle \FSurInv{(H^\perp)^{-1} s_\lambda}, s_\mu \rangle \\
            &= \langle s_\lambda, s_\mu[\omega(L_1) - \omega(L_2) + \omega(L_3) - \cdots] \cdot (1 - e_1 + e_2 - e_3 + \cdots)\rangle \\
            &= (-1)^{m - n} \langle s_{\lambda^T}, s_{\mu^T}[L_1 + L_2 + L_3 + \cdots] \cdot H \rangle.
        \end{align*}
    \end{enumerate}
\end{theorem}
Before proving \Cref{theorem:summary}, we restate a basic result of plethystic calculus.
\begin{lemma}[{Negation Rule, \cite[Theorem~6]{MR2765321}}]\label{lemma:negation}
    Let $f \in \Lambda$ and $g \in \overline{\Lambda}$. If~$f$ is homogeneous, then \[f[-g] = (-1)^{\deg f} (\omega(f))[g].\]
\end{lemma}
\begin{proof}[Proof of \Cref{theorem:summary}]
    \begin{enumerate}[label=(\alph*)]
        \item The first equality is true by definition. The second follows from \Cref{remark:halladjoint}.
        \item The first equality is true by definition. The second is demonstrated in the second proof of \Cref{proposition:fsur}.
        \item The first equality is true by definition.
        
        For the second, Cadogan showed in 1971 that $\omega(L_1) - \omega(L_2) + \omega(L_3) - \cdots$ is the plethystic inverse of $H_+$~\cite{MR0284377}. In part (b), we showed that $\Fa_{\Sur}$ is adjoint to plethysm by $H_+$. Hence, $\Fa_{\Sur}^{-1}$ is adjoint to plethysm by $\omega(L_1) - \omega(L_2) + \omega(L_3) - \cdots$, as desired.

        For the third, by \Cref{lemma:negation} and the associativity of plethysm, we have
        \begin{align*}
            & \langle s_\lambda, s_\mu[\omega(L_1) - \omega(L_2) + \omega(L_3) - \cdots] \rangle \\
            =& \langle s_\lambda, s_\mu[-(L_1 + L_2 + L_3 + \cdots)[-p_1]] \rangle \\
            =& \langle s_\lambda, (s_\mu[-(L_1 + L_2 + L_3 + \cdots)]) [-p_1] \rangle \\
            =& \langle s_\lambda[-p_1], s_\mu[-(L_1 + L_2 + L_3 + \cdots)] \rangle \\
            =& \langle (-1)^m \omega(s_\lambda)[p_1], (-1)^n \omega(s_\mu)[L_1 + L_2 + L_3 + \cdots] \rangle \\
            =& (-1)^{m - n} \langle s_{\lambda^T}, s_{\mu^T}[L_1 + L_2 + L_3 + \cdots] \rangle 
        \end{align*}
        as desired.
        \item By part (a), we have
        \begin{equation}\label{equation:amulambda}
            a^\mu_{\lambda}
            = \lim_{n \to \infty} \langle \F{s_\lambda}, s_{\mu^{(n)}} \rangle 
            = \lim_{n \to \infty} \langle \FSur{s_\lambda} \cdot H, s_{\mu^{(n)}} \rangle.
        \end{equation}
        Now, we claim that for any $f \in \Lambda$, we have
        \begin{equation}\label{equation:limit}
            \lim_{n \to \infty} \langle f \cdot H, s_{\mu^{(n)}} \rangle = \langle f, s_\mu \cdot H \rangle.
        \end{equation}
        By linearity, it is enough to show \eqref{equation:limit} when $f = s_\nu$ for some partition~$\nu$. In that case, by the Pieri rule,
        \[
            \lim_{n \to \infty} \langle f \cdot H, s_{\mu^{(n)}} \rangle = \lim_{n \to \infty} \begin{cases} 1 & \mbox{if $\mu^{(n)}/\nu$ is a horizontal strip;} \\ 0 & \mbox{otherwise.}\end{cases}
        \]
        For~$n$ sufficiently large, it is easy to see that $\mu^{(n)}/\nu$ is a horizontal strip if and only if $\nu / \mu$ is a horizontal strip. So the above limit is equal to
        \[
            \begin{cases} 1 & \mbox{if $\nu/\mu$ is a horizontal strip;} \\ 0 & \mbox{otherwise,}\end{cases}
        \]
        which is just $\langle f, s_\mu \cdot H \rangle$. Hence, \eqref{equation:limit} indeed holds. Substituting \eqref{equation:limit} into \eqref{equation:amulambda} with $f = \FSur{s_\lambda}$, we obtain \[a^\mu_\lambda = \langle \FSur{s_\lambda}, s_\mu \cdot H \rangle.\] The result now follows from the fact that $\Fa_{\Sur}$ is adjoint to plethysm by $H_+$.
        \item Define the \emph{stable Frobenius transform} $\mathcal{A} \colon \Lambda \to \Lambda$ by $\mathcal{A} \{f\} = H^\perp \FSur{f}$. Part (d) implies that the matrix of $\mathcal{A}$ in the Schur basis is $[a_\lambda^\mu]$. Therefore, the matrix of $\mathcal{A}^{-1}$ in the Schur basis is $[b_\lambda^\mu]$. This proves the first equality.
        
        The second equality then follows from the fact that $\Fa_{\Sur}^{-1}$ is adjoint to plethysm by $\omega(L_1) - \omega(L_2) + \omega(L_3) - \cdots$.
        
        For the third equality, we again use \Cref{lemma:negation}, together with the fact that plethysm by $-p_1$ is an isometry and a ring automorphism:
        \begin{align*}
            &\langle s_\lambda, s_\mu[\omega(L_1) - \omega(L_2) + \omega(L_3) - \cdots] \cdot (1 - e_1 + e_2 - e_3 + \cdots)\rangle \\
            =& \langle s_\lambda, s_\mu[-(L_1 + L_2 + L_3 + \cdots)][-p_1] \cdot (1 - e_1 + e_2 - e_3 + \cdots)\rangle \\
            =& \langle s_\lambda[-p_1], s_\mu[-(L_1 + L_2 + L_3 + \cdots)] \cdot (1 - e_1 + e_2 - e_3 + \cdots)[-p_1]\rangle \\
            =& \langle (-1)^m \omega(s_\lambda)[p_1], (-1)^n \omega(s_\mu)[L_1 + L_2 + L_3 + \cdots] \cdot H\rangle \\
            =& (-1)^{m - n} \langle s_{\lambda^T}, s_{\mu^T}[L_1 + L_2 + L_3 + \cdots] \cdot H\rangle,
        \end{align*}
        as desired. \qedhere
    \end{enumerate}
\end{proof}
\begin{thebibliography}{MMMMMM}
\bibitem[1]{MR2122332} Andrews, George E.; Eriksson, Kimmo Integer partitions, Cambridge University Press, Cambridge, 2004, x+141 pages \href{https://alco.centre-mersenne.org/articles/10.5802/alco.365/#MR2122332}{Publisher reference}.
\bibitem[2]{MR4055931} Assaf, Sami H.; Speyer, David E. Specht modules decompose as alternating sums of restrictions of Schur modules, Proc. Amer. Math. Soc., Volume 148 (2020) no. 3, pp. 1015-1029 \href{https://alco.centre-mersenne.org/articles/10.5802/alco.365/#MR4055931}{Publisher reference}.
\bibitem[3]{MR0284377} Cadogan, C. C. The Möbius function and connected graphs, J. Combinatorial Theory Ser. B, Volume 11 (1971), pp. 193-200 \href{https://alco.centre-mersenne.org/articles/10.5802/alco.365/#MR0284377}{Publisher reference}.
\bibitem[4]{MR0102539} Chen, K.-T.; Fox, R. H.; Lyndon, R. C. Free differential calculus. IV. The quotient groups of the lower central series, Ann. of Math. (2), Volume 68 (1958), pp. 81-95 \href{https://alco.centre-mersenne.org/articles/10.5802/alco.365/#MR0102539}{Publisher reference}.
\bibitem[5]{MR1464693} Fulton, William Young tableaux: with applications to representation theory and geometry, London Mathematical Society Student Texts, 35, Cambridge University Press, Cambridge, 1997, x+260 pages \href{https://alco.centre-mersenne.org/articles/10.5802/alco.365/#MR1464693}{Publisher reference}.
\bibitem[6]{MR1397498} Graham, Ronald L.; Knuth, Donald E.; Patashnik, Oren Concrete mathematics: a foundation for computer science, Addison-Wesley Publishing Company, Reading, MA, 1994, xiv+657 pages \href{https://alco.centre-mersenne.org/articles/10.5802/alco.365/#MR1397498}{Publisher reference}.
\bibitem[7]{MR0927763} Joyal, André Foncteurs analytiques et espèces de structures, Combinatoire énumérative (Montreal, Que., 1985/Quebec, Que., 1985) (Lecture Notes in Math.), Volume 1234, Springer, Berlin, 1986, pp. 126-159 \href{https://alco.centre-mersenne.org/articles/10.5802/alco.365/#MR0927763}{Publisher reference}.
\bibitem[8]{MR1576896} Littlewood, D. E. Group Characters and the Structure of Groups, Proc. London Math. Soc. (2), Volume 39 (1935) no. 2, pp. 150-199 \href{https://alco.centre-mersenne.org/articles/10.5802/alco.365/#MR1576896}{Publisher reference}.
\bibitem[9]{MR2765321} Loehr, Nicholas A.; Remmel, Jeffrey B. A computational and combinatorial exposé of plethystic calculus, J. Algebraic Combin., Volume 33 (2011) no. 2, pp. 163-198 \href{https://alco.centre-mersenne.org/articles/10.5802/alco.365/#MR2765321}{Publisher reference}.
\bibitem[10]{MR1475463} Lothaire, M. Combinatorics on words, Cambridge Mathematical Library, Cambridge University Press, Cambridge, 1997, xviii+238 pages \href{https://alco.centre-mersenne.org/articles/10.5802/alco.365/#MR1475463}{Publisher reference}.
\bibitem[11]{MR0067884} Lyndon, R. C. On Burnside’s problem. II, Trans. Amer. Math. Soc., Volume 78 (1955), pp. 329-332 \href{https://alco.centre-mersenne.org/articles/10.5802/alco.365/#MR0067884}{Publisher reference}.
\bibitem[12]{MR3443860} Macdonald, I. G. Symmetric functions and Hall polynomials, Oxford Classic Texts in the Physical Sciences, The Clarendon Press, Oxford University Press, New York, 2015, xii+475 pages \href{https://alco.centre-mersenne.org/articles/10.5802/alco.365/#MR3443860}{Publisher reference}.
\bibitem[13]{MR1132426} Méndez, M. Tensor species and symmetric functions, Proc. Nat. Acad. Sci. U.S.A., Volume 88 (1991) no. 21, pp. 9892-9894 \href{https://alco.centre-mersenne.org/articles/10.5802/alco.365/#MR1132426}{Publisher reference}.
\bibitem[14]{MR4311085} Narayanan, Sridhar P.; Paul, Digjoy; Prasad, Amritanshu; Srivastava, Shraddha Character polynomials and the restriction problem, Algebr. Comb., Volume 4 (2021) no. 4, pp. 703-722 \href{https://alco.centre-mersenne.org/articles/10.5802/alco.365/#MR4311085}{Publisher reference}.
\bibitem[15]{MR4356269} Narayanan, Sridhar P.; Paul, Digjoy; Prasad, Amritanshu; Srivastava, Shraddha Polynomial induction and the restriction problem, Indian J. Pure Appl. Math., Volume 52 (2021) no. 3, pp. 643-651 \href{https://alco.centre-mersenne.org/articles/10.5802/alco.365/#MR4356269}{Publisher reference}.
\bibitem[16]{MR4245137} Orellana, Rosa; Zabrocki, Mike A combinatorial model for the decomposition of multivariate polynomial rings as S n -modules, Electron. J. Combin., Volume 27 (2020) no. 3, Paper no. 3.24, 18 pages \href{https://alco.centre-mersenne.org/articles/10.5802/alco.365/#MR4245137}{Publisher reference}.
\bibitem[17]{MR4275829} Orellana, Rosa; Zabrocki, Mike The Hopf structure of symmetric group characters as symmetric functions, Algebr. Comb., Volume 4 (2021) no. 3, pp. 551-574 \href{https://alco.centre-mersenne.org/articles/10.5802/alco.365/#MR4275829}{Publisher reference}.
\bibitem[18]{MR4295089} Orellana, Rosa; Zabrocki, Mike Symmetric group characters as symmetric functions, Adv. Math., Volume 390 (2021), Paper no. 107943, 34 pages \href{https://alco.centre-mersenne.org/articles/10.5802/alco.365/#MR4295089}{Publisher reference}.
\bibitem[19]{MR2035110} Reutenauer, Christophe Free Lie algebras, Handbook of algebra, Vol. 3 (Handb. Algebr.), Volume 3, Elsevier/North-Holland, Amsterdam, 2003, pp. 887-903 \href{https://alco.centre-mersenne.org/articles/10.5802/alco.365/#MR2035110}{Publisher reference}.
\bibitem[20]{MR4127906} Ryba, Christopher Resolving irreducible ℂS n -modules by modules restricted from GL n (ℂ), Represent. Theory, Volume 24 (2020), pp. 229-234 \href{https://alco.centre-mersenne.org/articles/10.5802/alco.365/#MR4127906}{Publisher reference}.
\bibitem[21]{MR4621625} Stanley, Richard P. Enumerative combinatorics. Vol. 2, Cambridge Studies in Advanced Mathematics, 208, Cambridge University Press, Cambridge, 2024, xvi+783 pages \href{https://alco.centre-mersenne.org/articles/10.5802/alco.365/#MR4621625}{Publisher reference}.
\bibitem[22]{MR3641124} Wiltshire-Gordon, John D. Representation Theory of Combinatorial Categories, Ph. D. Thesis, University of Michigan (2016) (113 pages) \href{https://alco.centre-mersenne.org/articles/10.5802/alco.365/#MR3641124}{Publisher reference}.
\bibitem[23]{MR1817706} Zabrocki, Mike Vertex operators for standard bases of the symmetric functions, J. Algebraic Combin., Volume 13 (2001) no. 1, pp. 83-101 \href{https://alco.centre-mersenne.org/articles/10.5802/alco.365/#MR1817706}{Publisher reference}.
\end{thebibliography}
\end{document}
